% 大爆炸（综述）
% license CCBYSA3
% type Wiki

本文根据 CC-BY-SA 协议转载翻译自维基百科\href{https://en.wikipedia.org/wiki/Big_Bang}{相关文章}。

\begin{figure}[ht]
\centering
\includegraphics[width=8cm]{./figures/7a2c26b638cb807a.png}
\caption{宇宙膨胀的时间线：其中空间在每个阶段以圆形剖面示意表示。左侧展示了宇宙暴胀时期的剧烈膨胀；中间部分则表示宇宙膨胀的加速阶段（艺术示意图；时间与尺度均非按比例绘制）。} \label{fig_BBang_1}
\end{figure}
“大爆炸”（The Big Bang）是一种物理理论，用以描述宇宙如何从最初高密度、高温的状态中膨胀而来。基于大爆炸概念的各种宇宙学模型能够解释广泛的天文现象\(^\text{[1][2][3]}\)包括轻元素的丰度、宇宙微波背景辐射（CMB），以及宇宙的大尺度结构。宇宙的均匀性——即所谓的视界与平坦性问题——通过“宇宙暴胀”得到解释：那是一段发生在宇宙最早期的加速膨胀阶段。对宇宙膨胀速率的精确测量将最初奇点的时间推算为约 137.87 ± 0.02 亿年前，这被认为是宇宙的年龄。大量实证证据强烈支持这一“大爆炸事件”，该理论如今已被广泛接受。\(^\text{[4]}\)

根据已知物理定律，将这种宇宙膨胀过程反向外推，模型显示宇宙早期处于极端炽热且高密度的原始状态。物理学目前仍缺乏一种被广泛接受的理论来完整描述大爆炸最初阶段的条件。\(^\text{[5]}\)随着宇宙膨胀，它逐渐冷却到足以形成亚原子粒子，随后又形成原子。这些原初元素——主要是氢，辅以少量氦和锂——在暗物质的引力作用下聚集，最终形成早期的恒星与星系。对超新星红移的观测表明，宇宙的膨胀正在加速，这一现象被归因于一种被称为“暗能量”的概念。

“宇宙膨胀”的思想最早由物理学家亚历山大·弗里德曼（Alexander Friedmann）于 1922 年提出，他在数学上推导出了“弗里德曼方程”。\(^\text{[6][7][8][9]}\)宇宙膨胀的最早观测依据被称为“哈勃定律”（Hubble’s Law），由物理学家爱德文·哈勃（Edwin Hubble）于 1929 年发表，他发现星系正以与其距离成正比的速率远离地球。独立于弗里德曼的理论工作与哈勃的观测，物理学家乔治·勒梅特（Georges Lemaître）于 1931 年提出，宇宙起源于一个“原初原子”（primeval atom），从而引入了现代意义上的“大爆炸”概念。1964 年，人们发现了宇宙微波背景辐射。随后几年中，对该辐射的观测表明：它在天空各方向上高度均匀，且能量—强度曲线的形状都与“大爆炸模型”预测的早期高温高密度状态相一致。至 20 世纪 60 年代末，大多数宇宙学家已确信竞争性的“稳恒态宇宙”模型是错误的。\(^\text{[10]}\)

然而，观测宇宙中仍存在若干现象尚未能由“大爆炸模型”充分解释。其中包括：物质与反物质数量不对称（即重子不对称性）、星系周围暗物质的具体性质，以及“暗能量”的起源。\(^\text{[11]}\)
\subsection{模型的特征}
\subsubsection{基本假设}
大爆炸宇宙学模型依赖于三项主要假设：物理定律的普适性、宇宙学原理，以及物质可被视作理想流体。\(^\text{[12]}\)物理定律的普适性是相对论的基本原则之一。宇宙学原理指出，在大尺度上，宇宙是均匀且各向同性的——无论从哪个方向观察、位于何处，宇宙都呈现相同的总体结构。\(^\text{[13]}\)所谓“理想流体”，是指没有黏度的流体，其压力与密度成正比。\(^\text{[14]: 49 }\)

这些假设最初被视为公设，但后来科学家努力对其逐一加以验证。例如，第一个假设（物理定律的普适性）已通过观测得到检验：在宇宙大部分存在时期内，精细结构常数（fine-structure constant）的最大可能偏差量级仅为 $10^{-5}$。\(^\text{[15]}\)支撑这些模型的关键物理定律——广义相对论——已在太阳系与双星系统尺度上通过了严格的检验。\(^\text{[16][17]}\)宇宙学原理亦得到高度验证：通过对宇宙微波背景（CMB）温度的观测，其一致性水平达到 $10^{-5}$。截至 1995 年，对 CMB 视界尺度的测量表明，宇宙的非均匀性上限约为 10\%。\(^\text{[18]}\)
\subsubsection{膨胀预测}
宇宙学原理极大地简化了广义相对论方程，得出了用于描述宇宙几何结构的弗里德曼–勒梅特–罗伯逊–沃尔克度规（FLRW metric）；结合理想流体假设，可推导出弗里德曼方程，用于描述宇宙几何随时间的演化。在这一层次的描述中，唯一的自由参数是质量-能量密度：宇宙的几何形态与其膨胀行为均由密度直接决定。\(^\text{[19]: 73 }\)大爆炸宇宙学的所有主要特征都与这些结果密切相关。\(^\text{[14]: 49 }\)
\subsubsection{质量–能量密度}
\subs\begin{figure}[ht]
\centering
\includegraphics[width=8cm]{./figures/818e9d5df373b8f0.png}
\caption{} \label{fig_BBang_2}
\end{figure}
在大爆炸宇宙学中，质量–能量密度决定了宇宙的形状与演化。通过将天文观测结果与已知的热力学和粒子物理定律结合，宇宙学家推算出了宇宙生命期内密度的组成比例。在当前的宇宙中，可见物质（如恒星、行星等）所占的密度不到 5\%；暗物质约占 27\%；而暗能量则占据剩余的 68\%。\(^\text{[20]}\)
\subsubsection{视界}
大爆炸时空的一个重要特征是粒子视界的存在。由于宇宙的年龄有限，且光速也是有限的，因此可能存在一些过去的事件，其光尚未有足够的时间到达地球。这就为我们能观测到的最遥远天体设置了一个限制，即“过去视界”。  
相反，由于空间在不断膨胀，更遥远的天体正以越来越快的速度远离我们，因此我们今天发出的光可能永远无法“追上”那些极远的天体。这就定义了一个“未来视界”，它限制了我们未来能够影响到的事件范围。  
这两种视界的存在都取决于描述宇宙膨胀的弗里德曼–勒梅特–罗伯逊–沃尔克（FLRW）度规的具体形式。\(^\text{[21]}\)

根据我们对早期宇宙的理解，可以推知宇宙存在“过去视界”；但在实践中，我们的观测受到早期宇宙不透明性的限制，因此视野无法在时间上无限向后延伸，尽管该视界在空间上不断退后。若宇宙的膨胀继续加速，还将存在“未来视界”。\(^\text{[21]}\)
\subsubsection{热化过程}
早期宇宙中的某些物理过程相对于宇宙膨胀速率而言过于缓慢，因而无法达到近似热力学平衡；另一些过程则足够迅速，可以实现热化。  
通常用来判断早期宇宙中的某一过程是否达到热平衡的参数，是该过程速率（通常为粒子碰撞频率）与哈勃参数的比值。比值越大，意味着粒子在彼此分离之前有更多时间完成热化过程。\(^\text{[22]}\)

宇宙能量密度各组成部分的相对分布估计图。（2015 年 2 月，欧洲主导的普朗克宇宙学探测计划研究团队发布了新的精确数据：普通物质占 4.9\%，暗物质占 25.9\%，暗能量占 69.1\%。）
\subsection{时间线}
根据大爆炸模型，宇宙在最初阶段极其炽热而致密，此后一直在膨胀与冷却之中。
\subsubsection{奇点}
现有的物理理论无法描述大爆炸发生的确切瞬间。\(^\text{[5]}\)如果仅用经典广义相对论将宇宙的膨胀反推至过去，会得到一个具有无限密度与温度的引力奇点，该奇点存在于有限的过去时间中。\(^\text{[23]}\)然而，这种经典引力理论被认为不足以描述在如此极端条件下的物理。\(^\text{[19]: 275}\) 因此，这个奇点在大爆炸语境中的真实含义仍不明确。\(^\text{[24]}\)

广义相对论所能适用的最早时刻称为普朗克时间（Planck time）。\(^\text{[19]: 274}\) 在此之前，即所谓的普朗克纪元（Planck epoch），宇宙温度接近普朗克温标（约 $10^{32}$ K 或 $10^{28}$ eV），量子引力效应被认为占主导地位。  
迄今为止，人类尚未建立一个被普遍接受的量子引力理论；在高于普朗克能标的区域，未知的物理规律可能会影响宇宙的膨胀历史。
\subsubsection{暴胀与重子生成}
由于缺乏可观测数据，大爆炸最早阶段的研究仍主要基于推测。在最常见的模型中，宇宙在最初阶段具有极高的能量密度、温度与压力，并以极快的速度均匀膨胀与冷却。  
在膨胀约 $10^{-43}$ 秒之前的普朗克纪元中，四种基本相互作用——电磁力、强核力、弱核力与引力——统一为一种力。\(^\text{[25]}\)此时，宇宙的特征尺度长度为普朗克长度（$1.6\times10^{-35}$ 米），温度约为 $10^{32}$ 摄氏度；在这种条件下，“粒子”的概念本身已失去意义。对这一时期的准确理解有待量子引力理论的发展。\(^\text{[26][27]}\)普朗克纪元之后，在约 $10^{-43}$ 秒时进入\emph{大统一纪元}，此时引力从其他三种力中分离出来，因为宇宙温度继续下降。\(^\text{[25]}\)

大约在膨胀 $10^{-37}$ 秒时，一次相变引发了宇宙暴胀（cosmic inflation），期间宇宙以指数级速度膨胀，不受光速不变性的约束，温度下降了约十万倍。这一思想的提出是为了解释平坦性问题（flatness problem）——即宇宙的物质与能量密度非常接近产生平坦几何所需的临界密度。换言之，宇宙整体没有几何曲率。由于海森堡不确定性原理产生的微观量子涨落在暴胀期间被“冻结”，并被放大为后来形成宇宙大尺度结构的种子。\(^\text{[28]}\)在约 $10^{-36}$ 秒时进入电弱纪元（electroweak epoch），此时强核力从其他力中分离，只剩电磁力与弱核力保持统一。\(^\text{[29]}\)

当前可观测到的所有星系的质量–能量最初都集中在一个半径约 $4\times10^{-29}$ 米的球体中，到暴胀结束时该球体的半径约增至 0.9 米。\(^\text{[30]: 202}\) 随后进入再加热阶段（reheating），暴胀子场（inflaton field）衰变，使宇宙温度升高到足以产生夸克–胶子等离子体及各种基本粒子的水平。\(^\text{[31][32]}\)当时温度极高，粒子的随机运动达到相对论速度，各类粒子与反粒子对在碰撞中不断产生与湮灭。\(^\text{[33]}\)在某个阶段，发生了一种未知的反应称为重子生成（baryogenesis），它违反了重子数守恒定律，使夸克与轻子相对于反夸克与反轻子产生极微小的不对称，约为三千万分之一。这种微小的不对称导致了当今宇宙中物质相对于反物质的主导地位。\(^\text{[34]}\)



